%%% FIELD proof-first-lossy-size
By the declared type of \(n\), the admissible population sizes are the positive even integers.  Consequently,
\[
 \{n\in 2\mathbb Z_{>0}:n\leq 8\}=\{2,4,6,8\}.
\]
Theorem \(\mathrm{thm{:}n24\mbox{-}equality}\) gives
\[
 \rho_2^{\mathrm{CR}}=\rho_2^W=\frac14
 \quad\text{and}\quad
 \rho_4^{\mathrm{CR}}=\rho_4^W=1-\frac{\sqrt3}{2}.
\]
Theorem \(\mathrm{thm{:}n6\mbox{-}seven\mbox{-}state\mbox{-}saddle}\) gives
\[
 \rho_6^{\mathrm{CR}}=\rho_6^W=s_6^2.
\]
Finally, Theorem \(\mathrm{thm{:}n8\mbox{-}strict\mbox{-}separation}\) gives the strict inequality
\[
 \rho_8^{\mathrm{CR}}<\rho_8^W.
\]
Thus equality holds at each admissible size below eight and fails at eight.  It follows that
\[
 \{n\in 2\mathbb Z_{>0}:n\leq8,\ \rho_n^{\mathrm{CR}}=\rho_n^W\}
 =\{2,4,6\}.
\]
Because \(W=2(X_1-X_0)\) is precisely the stipulated compression of the full observation \((X_1,X_0)\), the strict inequality at \(n=8\) says exactly that this compression is minimax-lossy there.  Since no other admissible positive even size lies between six and eight, \(n=8\) is the first such size.
